\documentclass[termpaper,assignment]{pmk-graduate}

\begin{document}
% Учебный год
\year{2025-2026}

% Дисциплина
\subject{Системное программирование}

% Студент
\student{Митрофанов Денис Алексеевич
\brief Д.\,А.\,Митрофанов
\group 37
\english Mitrofanov Denis Alekseevich
}

% Научный руководитель
\superviser{Дудаков Сергей Михайлович
\brief С.\,М.\,Дудаков
\position декан факультета ПМиК
\degree д.\,ф.-м.\,н. % Учёная степень. При отсутствии убрать строку
\rank профессор % Учёное звание. При отсутствии убрать строку
}

% Образовательная программа. См. перечень ниже
\program{ics}
% Бакалавриат:
% ics - Информатика и компьютерные науки
% swi - Инженерия программного обеспечения
% swiai - Программная инженерия в искусственном интеллекте
% aida - Искусственный интеллект и анализ данных
% mm - Математическое моделирование
% sa - Системный анализ
% icpim - Инженерное компьютерное проектирование и индустриальная математика
% acse - Прикладная информатика в экономике
% acsm - Прикладная информатика в мехатронике
% icmrs - Интеллектуальное управление в мехатронных и робототехнических системах
% Магистратура:
% macsae - Прикладная информатика в аналитической экономике
% mista - Интеллектуальные системы. Теория и приложения
% memsp - Системное программирование
% memsa - Системный анализ
% mitcd - Информационные технологии в управлении и принятии решений


% Заголовок работы
\title{Интеллектуальные методы корректировки ограничения целостности по образцу
    \english Intelligent methods for correcting integrity constraints based on a sample}

% Цель работы
\goal{Разработка системы для коррекции ограничений целостности по образцу с использованием интеллектуальных методов.}

% Задачи для разработки в работе
\begin{tasks}
  \task{Выбрать метод, подходящий для достижения поставленной цели и обосновать его выбор.
  }
  \task{Описать теоретические основы выбранного метода.
  }
  \task{Реализовать предложенный метод программно для коррекции ограничения целостности по образцу.
  }
  \task{Расписать и объяснить используемые в программе механизмы.
  }
  \task{Провести тестирование и анализ реализованной программы на практике.
  }
\end{tasks}

% Команда для генерации задания на работу
\makeassignment
\end{document}
